\section{결과}
\label{sec:results}

결과는 연구 질문의 순서로, 해당 그림과 표를 따라 보고한다.
모든 차이는 12개 코호트의 확률효과 통합 추정치와 95\% 신뢰구간이며, 괄호 안의 분수는 통합 추정치와 부호가 같은 코호트의 수이다(12/12는 모든 코호트에서 방향이 같다는 뜻이다).
AUC는 무작위로 고른 탈락 팀과 생존 팀의 쌍에서 모델이 탈락 팀에 더 높은 위험을 줄 확률이므로, AUC 차이 0.01은 이 확률이 1\%p 높다는 뜻이다.
차이의 95\% 신뢰구간이 0보다 크고 점추정치가 0.005 이상이면 개선으로, 90\% 신뢰구간이 $\pm0.005$ 안에 있으면 실질적 동등, 즉 차이가 있더라도 0.5\%p 미만으로 읽는다.

\begin{figure}[t]
\centering
\includegraphics[width=\columnwidth]{fig_rq1_effects.pdf}
\caption{단계 구간별 서식지 맥락의 효과}
\label{fig:rq1}
\end{figure}

\begin{table*}[t]
\centering
\caption{설계별 예측 성능과 원래 과제 성능}
\label{tab:main}
\footnotesize
\setlength{\tabcolsep}{4pt}
\renewcommand{\arraystretch}{1.15}
\begin{tabular}{lcccccccc}
\hline
 & \multicolumn{6}{c}{본 과제: 다음 단계 탈락} & \multicolumn{2}{c}{KRAFTON의 원래 과제} \\
\cline{2-7}\cline{8-9}
설계 & AUC & 단계 안 AUC & 경기--단계 내부 AUC & 로그 손실 & Brier & 단계별 보정 오차 & 우승 팀 정확도 & C-index \\
\hline
본 연구, GBDT & 0.778 & 0.752 & \textbf{0.757} & \textbf{0.489} & \textbf{0.160} & \textbf{0.013} & 0.351 & 0.654 \\
본 연구, PTR-Former & \textbf{0.779} & \textbf{0.757} & \textbf{0.757} & 0.490 & \textbf{0.160} & 0.017 & 0.386 & 0.663 \\
KRAFTON 적용판 & 0.741 & 0.741 & 0.748 & 0.518 & 0.171 & 0.064 & \textbf{0.417} & \textbf{0.665} \\
KRAFTON + 단계 기저위험 & 0.763 & 0.741 & 0.748 & 0.504 & 0.165 & 0.017 & \textbf{0.417} & \textbf{0.665} \\
Sagong 적용판(XGBoost) & 0.686 & 0.632 & 0.645 & 0.562 & 0.188 & 0.016 & --- & --- \\
\hline
\multicolumn{9}{l}{\emph{같은 KRAFTON 구조에서 학습 목적만 바꾼 경우}} \\
Cox 목적 & 0.742 & 0.731 & 0.745 & 0.517 & 0.171 & 0.056 & \textbf{0.387} & \textbf{0.661} \\
다음 단계 교차엔트로피 목적 & \textbf{0.776} & \textbf{0.755} & \textbf{0.753} & \textbf{0.491} & \textbf{0.161} & \textbf{0.016} & 0.362 & 0.653 \\
\hline
\end{tabular}
\\[2pt]
\parbox{\textwidth}{\scriptsize 12개 코호트 평균. 굵은 글씨는 각 블록에서 가장 좋은 값이다. AUC와 정확도, C-index는 클수록, 로그 손실, Brier, 보정 오차는 작을수록 좋다. 우승 팀 정확도와 C-index는 KRAFTON 구현이 정의한 지표이며, 본 연구의 PTR-Former는 순위 출력을 사용하였다.}
\end{table*}

\begin{figure*}[t]
\centering
\includegraphics[width=\textwidth]{fig_contrast.pdf}
\caption{단계별 보정과 설계 요인의 분해}
\label{fig:contrast}
\end{figure*}

%======================================================================
\subsection{RQ1: 서식지 맥락과 다음 단계 탈락 위험}
\label{sec:results:rq1}

그림~\ref{fig:rq1}의 계수는 각 요인이 한 표준편차 늘 때 탈락 위험의 로그 오즈가 바뀌는 양으로, 0보다 크면 위험을 높이고 작으면 낮추며, 지수를 취하면 탈락 오즈가 몇 배가 되는지를 나타낸다.
모든 단계를 합친 모형에서 주변 경쟁 밀도는 $+0.293$ $[0.242, 0.343]$(12/12)로, 주변의 적 팀이 한 표준편차 많으면 탈락 오즈가 약 1.34배가 되었다.
이동 부족분은 $+0.067$ $[0.051, 0.083]$(12/12, 약 1.07배), 남은 팀원 수는 $-0.632$ $[-0.713, -0.551]$(11/11, 약 0.53배)였다.
세 요인 모두 신뢰구간이 0을 포함하지 않고 평가 가능한 모든 코호트에서 방향이 같으므로, 경쟁과 이동 부담은 위험을 높이고 남은 팀원은 위험을 낮춘다고 해석된다.
반면 구역 밖 깊이는 방향이 코호트마다 달라(8/12), 상관이 높은 이동 부족분과 구분되는 효과를 확인할 수 없었다.

효과의 크기는 단계에 따라 달랐다.
구역 안쪽 깊이의 계수는 초반 $-0.028$(신뢰구간이 0을 포함, 7/12)에서 중반 $-0.289$, 후반 $-0.410$(모두 12/12)으로 커졌는데, 이는 구역이 작아질수록 경계에서 깊이 들어가 있는 것이 더 큰 보호가 된다는 뜻이다.
경쟁 밀도의 계수는 초반 $+0.376$에서 후반 $+0.229$로 작아졌고, 구간 간 차이는 우연으로 보기 어려웠다($p=5.9\times10^{-5}$).
확률 단위로 보면, 주변의 적 팀이 하나 늘 때 다음 단계 탈락확률은 평균적으로 초반 8.0\%p, 중반 4.1\%p, 후반 3.8\%p 높아졌고, 팀원이 한 명 더 남아 있으면 각각 8.5\%p, 12.6\%p, 14.1\%p 낮아졌다.

이 요인들은 적은 수의 변수로도 위험의 상당 부분을 설명한다.
필요 속도, 주변 경쟁 밀도, 생존 팀 수의 로그, 남은 팀원 수의 네 변수만 쓴 간이 모형은 GBDT가 우연 수준(AUC 0.5)을 넘어 얻은 판별력의 78.8\% $[77.2, 80.3]$를 얻었다.
즉 150개 특징을 쓰는 모델이 가려내는 것의 약 5분의 4를 해석 가능한 네 변수로 설명할 수 있다.

%======================================================================
\subsection{RQ2: 기존 설계와의 비교}
\label{sec:results:rq2}

표~\ref{tab:main}에서 본 연구의 설계는 AUC 0.778(GBDT)과 0.779(PTR-Former)로, KRAFTON 적용판(0.741), 단계 기저위험을 결합한 KRAFTON(0.763), Sagong 적용판(0.686)보다 높았다.
PTR-Former와 기저위험을 결합한 KRAFTON의 차이는 $+0.016$ $[0.014, 0.019]$로, 탈락 팀과 생존 팀의 쌍을 올바르게 줄 세울 확률이 1.6\%p 높고 12/12 코호트에서 같은 방향이어서 개선으로 판정된다.
Sagong 적용판과의 차이는 $+0.089$ $[0.082, 0.096]$(12/12)로 더 컸다.

그림~\ref{fig:contrast}(a)는 단계마다 실제 탈락률(검은 선)과 각 설계의 평균 예측확률을 겹쳐 보이며, 두 선이 겹칠수록 그 설계가 단계별 위험 수준을 바르게 나타낸다는 뜻이다.
실제 탈락률은 단계 1의 19.5\%에서 단계 7의 50.6\%로 오르지만, KRAFTON 적용판의 예측은 26.4\%에서 29.8\%로 거의 평평하여 초반에는 위험을 과대평가하고 단계 7에서는 실제 50.6\%를 29.8\%로 과소평가한다.
단계별 보정 오차는 KRAFTON이 0.064로 단계별 평균 예측이 실제와 평균 6.4\%p 어긋난 반면, 본 연구의 GBDT는 0.013(1.3\%p)이었고, KRAFTON에 단계 기저위험을 결합하면 0.017로 줄었다.
이는 상대위험 목적이 단계별 위험 수준을 담지 못한다는 \ref{sec:method:rq2}절의 분석과 일치한다.

이 차이가 단계 수준만 맞힌 결과는 아니다.
같은 단계의 표본끼리만 비교한 단계 안 AUC에서 본 연구의 설계가 $+0.017$ $[0.013, 0.021]$(10/12), 같은 경기의 같은 순간에 살아 있던 팀끼리만 비교한 경기--단계 내부 AUC에서도 $+0.010$ $[0.006, 0.015]$(11/12) 높았다.
즉 같은 상황에 놓인 팀들 사이에서도 누가 먼저 탈락할지를 더 잘 가렸다.

표~\ref{tab:main}의 오른쪽 두 열은 KRAFTON의 원래 과제 성능이다.
우승 팀 정확도는 각 시점에 살아 있는 팀 가운데 최종 우승 팀을 1위로 맞힌 비율이고(무작위로 고를 때 0.182), C-index는 탈락 순서를 맞힌 팀 쌍의 비율이다.
이 과제에서는 KRAFTON 적용판이 0.417과 0.665로 가장 높았고 본 연구의 GBDT는 0.351과 0.654였으며, 같은 KRAFTON 구조에서도 Cox 목적은 원래 과제에서, 다음 단계 교차엔트로피 목적은 본 과제에서 더 높았다.
따라서 적용판은 원래 과제에 맞게 학습되었고, 각 학습 목적은 서로 다른 질문에 강하다.

%======================================================================
\subsection{RQ3: 설계 요인의 분해}
\label{sec:results:rq3}

그림~\ref{fig:contrast}(b)와 (c)의 막대 하나는 요인 하나를 바꿀 때 AUC가 변한 양이며, 진한 막대가 큰 변화이다.
Sagong 경로(b)에서는 같은 GBDT에 서식지 맥락 블록만 더해도 AUC가 0.699에서 $+0.047$ 올랐고, 나머지 입력을 본 연구의 나머지 팀 특징으로 바꾸면 $+0.029$가 더해졌다.
전체 차이의 약 60\%가 서식지 맥락 블록에서 나왔으므로, Sagong 설계와의 차이는 주로 입력 표현, 특히 구역에 대한 팀의 처지를 담았는지에서 온다.

KRAFTON 경로(c)에서는 입력과 학습 표본을 본 연구에 맞추는 단계가 $+0.001$로 거의 변화가 없었다.
반면 Cox 점수에 단계별 기저위험을 결합하면 $+0.016$ $[0.010, 0.022]$, 학습 목적을 다음 단계 교차엔트로피로 바꾸면 $+0.018$ $[0.016, 0.020]$이 더해졌다(모두 12/12).
목적 교체의 이득은 단계 안 AUC($+0.022$)와 경기--단계 내부 AUC($+0.009$)에서도 남아, 단계 수준을 맞추는 효과에 그치지 않았다.
명시적 단계 절편($-0.001$), 아키텍처($+0.001$), 단계 축과 관계 편향($+0.002$)은 모두 실질적 동등이었다.
따라서 전체 차이 0.037 가운데 약 90\%(0.034)가 단계별 기저위험과 학습 목적, 곧 문제를 어떻게 정식화했는지에서 나왔다.

사전 지정한 다섯 구조 비교도 모두 실질적 동등이었고, REML과 Hartung--Knapp 방식으로 다시 통합해도 판정은 같았다(부록).
팀 축의 효과는 $+0.0027$ $[0.0015, 0.0038]$(Holm $p<0.001$)로 통계적으로는 0보다 컸지만 크기가 0.3\%p에 그쳐 실질적 이득으로 보기 어렵다.
학습 예산을 두 배로 늘려도 신경망의 성능은 오르지 않았고(10개 실행 중 9개에서 $-0.0004$에서 $-0.0027$), 순위 보조 손실의 효과도 0.001 이하였으므로, 구조의 기여가 작은 것을 학습 부족 탓으로 보기는 어렵다.

입력끼리 비교하면 본 연구의 입력과 재구성한 KRAFTON 입력은 같은 GBDT에서 차이가 $+0.001$ $[-0.002, 0.004]$로 실질적으로 같았다.
KRAFTON 입력에도 구역까지의 거리와 주변 팀 수가 들어 있어(표~\ref{tab:blocks}) 두 입력의 정보는 상당 부분 겹칠 수 있으나, 둘을 합치면 $+0.014$ $[0.013, 0.016]$(12/12) 높아져 서로 다른 정보도 담고 있었다.
가장 큰 두 코호트의 학습 곡선에서는 학습 경기를 모두 써도 AUC가 계속 올랐고, PTR-Former는 전체 학습 경기를 쓸 때에만 GBDT를 넘어섰다(부록).
